\chapter{Conclusion}
\label{chap:conclusion}

\section{What was asked and what was built}
\label{sec:concl-summary}

The always-on cognitive architecture of \citet{Pasquali2025AlwaysOn,Pasquali2026ISEED} could perceive continuously, segment experience into episodes and recognise the social character of a room, yet nothing in it decided that any of this mattered. Its behaviour module had no interface through which an internal signal could act. This thesis built one, an embodied behaviour architecture in which such signals modulate eligibility, executive policy and dialogue, and installed a first value on it: a synthetic orexigenic variable that depletes over time and with each social action, and whose distance from its set point makes an event matter to the robot. Together the two close, in minimal form, the loop the programme had left open.

\section{Verdict on the research questions}
\label{sec:concl-questions}

The drive performed the functions classical homeostasis requires (RQ1): monitoring ran whether or not anyone interacted with the robot, deficit detection was exact at both thresholds, remote pings fired only from deficit states, and interaction progression reorganised once the critical threshold was crossed. Feeding was more frequent under deficit in both phases, but recovery was not reliable after every request (RQ2). Caregivers scaled meal size to the deficit and every observed entry into Starving was resolved through feeding, yet \CoreStarvingOnsetUnfed{} of the \CoreStarvingInteractions{} interactions beginning in critical deficit went unfed, only \CoreTelegramMessages{} of \CoreTelegramPings{} remote notifications drew a reply, and the longest depletions came when nobody was present. Regulatory memory encoded attributed outcomes exactly and expressed them at the eligibility gate, lowering the threshold for caregivers who had fed the robot and raising it in the one case of negative affinity (RQ3). The drive made the robot's need actionable, but the ratings give no evidence that it made the need credible (RQ4). Mean hunger ratings were below the midpoint for naturalness and genuineness and above it for helping motivation, although none differed significantly from the midpoint after correction. Helping motivation exceeded naturalness in an exploratory paired comparison. Written accounts identified repetitive, poorly timed requests and an unchanging non-verbal manner. Participants still found the robot likeable and benign, and \CoreCustomPartOfEnvironmentAgreement{} of \CoreCohortN{} felt it belonged to the office.

\section{Contribution and broader impact}
\label{sec:concl-contribution}

The architectural contribution is a behaviour layer through which an internal value can act inside the existing always-on loop without disturbing the substrate beneath it, together with verification, through the orexigenic drive, that it works as specified. What generalises is the gating pathway: a thresholded scalar variable can set executive policy directly and, through the affinity it produces, shift eligibility for later approaches.

The empirical contribution is that effectiveness and credibility, measured separately in the same deployment, diverged. A drive can enter control cleanly, produce the regulatory behaviour it was designed to produce, and still be disbelieved by the people whose cooperation it depends on. Studies reporting only recovery statistics, or only perceived animacy, would each have missed one of the two findings reported here. Work on artificial motivation should therefore treat regulatory effectiveness and interpretive credibility as separate endpoints and stop assuming that a correctly implemented drive will be read as genuine.

For robots that must ask for help, the absence of hunger expression in face, prosody and head motion is one plausible contributor to this credibility gap (Section~\ref{sec:disc-central-finding}). The guilt-inducing Telegram wording provides another explanation; both require controlled evaluation.

For social robotics more broadly, the deployment gives the autonomy argument of Section~\ref{sec:intro-homeostasis-drive} empirical weight. Caregivers sized their help to the robot's need, yet objected to requests that arrived at the wrong moment. Asking well, in proportion and in time, is part of what such a robot's autonomy consists of.

\section{Future work}
\label{sec:concl-future-work}

\citet{Pasquali2026ISEED} set two next steps for iSEED. The first is structural. Perception, segmentation, memory and recognition still run one after another as a pipeline, which was built to validate how experience is organised. Their framework of coupled dynamical loops needs modules with their own temporal dynamics: perception should stay locked to the changing world, while memory and motivation follow slower profiles of consolidation, reinforcement, decay and shifting salience. The orexigenic drive already has such a profile, a reserve that drains over hours, but it reads the memory only through a context label and never changes how memory consolidates or decays; coupling the two on their own timescales is the natural extension. The second step is to close the cognitive loops, so the agent decides how to act from the perceived context, its past experience and its values. This thesis closes one such loop with a single value and a designed policy. The policy is not yet learned, and the value does not yet reach memory.

The open problems \citet{Pasquali2025AlwaysOn} listed (Section~\ref{sec:icub-ecosystem}) are partly answered here. The value system now exists, with energy, one of the drives they name, and the online recognition path labels an episode before it ends. What remains concerns memory: segmentation that fragments busy periods (they propose a short-term window of about 30~s), consolidation without forgetting or reinforcement, and perception without semantics.

Several of these gaps share one cause, the absence of a learned model of how the scene evolves. A joint-embedding predictive architecture (JEPA) could supply one. V-JEPA~2.1 \citep{MurLabadia2026VJEPA21} learns from unlabelled images and video by predicting masked content in a learned representation space instead of in pixels, and its features stay spatially structured and temporally consistent. They support action anticipation and, with trained action-conditioned models, robot grasping and navigation. \citet{Pasquali2026ISEED} want an ``evolutionary background'' to bootstrap development, and such an encoder could play that role without replacing development. Its prediction error could do three jobs. It would replace the three-standard-deviation boundary rule with a learned segmentation signal. It would turn temporal expectation, now a lookup of which context usually occurs at a given hour, into a prediction that live observation can contradict. And it would give a forgetting policy the salience cue it needs. Because stored episodes keep their raw frames and audio, each role can first be tested on the existing office and conference data.

Learned action selection is the matching step on the behaviour side. A transformer-based policy can treat control as sequence modelling: it predicts the next action from past states and actions, given a target return. The deployment logs already hold such sequences (context label, hunger state, candidate person, action and resource outcome), with drive reduction as the natural return. A model trained on them could replace the hand-set priority weights and the designed policy, closing the loop in the learned form both iSEED papers aim at, while the fixed escalation to every subscriber stays as a safety floor. The same learning could set the drive's own set point and action costs from experience, which would complete the causal grounding that iSEED asks for and this thesis met only in part.

The limitations of Section~\ref{sec:disc-identification} set the most immediate work. Each points to a specific fix:
\begin{itemize}
    \item \emph{Inexpressive drive.} Hunger should reach face, prosody and head motion, and requests should vary with severity instead of repeating the same wording.
    \item \emph{No anticipation.} Temporal expectation could make the drive allostatic, forecasting depletion against who is likely to be present so the robot asks before the room empties.
    \item \emph{Missing controls and confounded phases.} Days with the drive switched off but action costs and resource logging kept, a neutral-affinity condition, and a fresh cohort or counterbalanced phase order would separate hunger expression from carried-over roles, time of day and caregiver availability.
    \item \emph{Uncertain attribution.} Merging identity spelling variants at enrolment, binding each meal to the person who presented the card, and matching Telegram replies by delivery receipt and subsequent feeding would replace timing-based rules with verified links.
    \item \emph{Narrow sample.} Longer deployments at several sites, with robot-naive participants and questionnaires repeated over weeks, would separate novelty from settled acceptance \citep{DeGraaf2016LongTermAcceptance}.\end{itemize}
A second drive would then test competition between needs \citep{Rosado2022DriveCompetition}.
